% LNCS appendix fragment. Include this file after \appendix with
% \input{appendixC}. It uses amsmath and amssymb and is independent of an
% algorithm package.

\section{Canonicalization Algorithm and Rewrite Theory}
\label{app:canon-rules}

This appendix makes explicit the proposed canonicalizer and its abstract
rewrite contract.  The Lean development proves finite-orbit normalization and
the Boolean, list, finite-function, and relation equations identified below.
Figure~\ref{alg:canon-g} and the Alloy/Java names are a refinement target for
those abstract objects, not a proved implementation correspondence.  An arrow
below denotes a proposed representative orientation unless the text marks the
operation as structural cleanup.  Publication run \path{db9f89bf-0965-4d74-8080-d9191d5f1aec}
records rule set \path{canonical-equivalences-v3-explicit-laws}.  This binds the
measured configuration but does not prove that Java realizes this appendix's
abstract rules or contracts.

\subsection{Notation and Preconditions}

Let $\mathcal G=(U,M,H)$ be a flexible-arity typed slotted e-graph.  The parent
assignment $U$ maps each e-class identifier to an immediate parent invocation,
with a root mapped to its identity invocation;
$\operatorname{find}_\Gamma$ composes these links to return a leader invocation
in caller context $\Gamma$.  The map
$M(a)=(\tau_a,S_a,B_a,G_a)$ stores the result type, exposed typed slots,
canonical shapes, and justified slot-symmetry group of e-class $a$, and $H$ is
the shape collision index.  A typed e-node has the form
\[
  n=f_\theta(q_1,\ldots,q_r),
  \qquad
  \Sigma(f)=\forall\bar\alpha.(\kappa_1,\ldots,\kappa_r)\Rightarrow\tau,
\]
where $q_i\in\operatorname{Port}_\Gamma(\theta(\kappa_i))$.  The schemas
$\texttt{Seq}$, $\texttt{Bag}$, and $\texttt{Set}$ respectively preserve order
and multiplicity, quotient by permutation while preserving multiplicity, and
quotient by both permutation and idempotence.  A first-class
$\texttt{BindBlock}(\beta,\kappa)$ carries a complete descriptor $\beta$, a
fresh occurrence-local bound context, and precisely the certified finite group
$\operatorname{Aut}(\beta)$.  Fix a total structural order $\triangleleft$ over
complete well-typed port and node keys, and a total order $\prec$ over completed
typed shapes.  The latter compares the operator, type instantiation, atom, slot
interface, and normalized ports.  Fix also the same total map order as in
the main text,
$\triangleleft_{\mathrm{map}}:=\prec_{\mathrm{ren}}$, for deterministic witness selection, and
let $\triangleleft_{\mathrm{lex}}$ compare a shape--map pair first by
$\prec$ and then by $\triangleleft_{\mathrm{map}}$.

At the specification level, leader normalization may reduce support, so the
proposed graph operation distinguishes the input context from the effective
kernel context.  For a node $n$ with
$\Gamma_0=\operatorname{slots}(n)$, its intended structural extraction is
\begin{align*}
\operatorname{leaderKernelTrace}_{\mathcal G}(n)
  &=(K_{\mathcal G}(n),\iota_n,\xi_n),\\
\Delta_n&=\operatorname{slots}(K_{\mathcal G}(n))\subseteq\Gamma_0,
&\iota_n&:\Delta_n\hookrightarrow\Gamma_0.
\end{align*}
Here $K_{\mathcal G}(n)$ and $\xi_n$ denote the proposed kernel and trace.
\leanref{TSG-APP-001}{TypedSlottedEGraphsPaper.ProfileRules.tsgApp001_kernelTraceReplayContract}
The exact formal claim \textsc{app-001} uses an indexed trace interface.  For an
indexed family $\mathsf{Term}:\mathbb N\to\mathsf{Type}$, an extractor returns
an effective size $k$, a kernel in $\mathsf{Term}(k)$, and a structural trace
with dependent endpoints
$\langle n,\mathit{input}\rangle$ and $\langle k,\mathit{kernel}\rangle$.
For any two trace interpreters, each supplying an equality for every trace
label, replay proves equality of the corresponding endpoint
denotations.  This theorem neither identifies those equalities with the
paper's $\operatorname{EqCert}$ type nor proves that Java extraction emits the
required trace; those are separate refinement obligations.

For a finite typed context $\Delta$, $\operatorname{Can}(\Delta)$ is the
context containing the same number of
slots of each type, named from a fixed type-indexed alphabet.  The shape
witness for $n$ is
\[
  \sigma_n\in\operatorname{TRen}(\operatorname{Can}(\Delta_n),\Delta_n),
\]
which maps canonical slots bijectively to the effective source support.  Its
transport into the original ambient context is the typed embedding
\[
  \omega_n=\iota_n\circ\sigma_n
  \in\operatorname{TEmb}(\operatorname{Can}(\Delta_n),\Gamma_0).
\]
The map $\omega_n$ is classified as a typed embedding; it is a renaming in the
support-preserving case $\Delta_n=\Gamma_0$.  This separation retains the
shape's exact slots, its effective support, and its original ambient support,
including coordinates removed by find.  Binder-block automorphisms act on
exactly the bound coordinates admitted by the complete
descriptor: domain, quantifier, cardinality/multiplicity, disjointness class,
dependency constraints, and type.  Each block occurrence is alpha-aligned with
a fresh canonical variant of the descriptor's bound context; the descriptor
group is transported to that occurrence context before it acts.  This local
action is disjoint from the one global free-slot renaming.

As the smallest strict-support case, suppose
$U(a)=m*\ell$ with a proper embedding $m:S_\ell\hookrightarrow S_a$.
Then $\operatorname{wrap}(\operatorname{id}_{S_a}*a)$ has input support $S_a$,
whereas its leader kernel $\operatorname{wrap}(m*\ell)$ has support
$m[S_\ell]$.  The construction below therefore uses
$\operatorname{Can}(m[S_\ell])$ and keeps
the inclusion into $S_a$ separately.  Its shape witness is the bijection onto
the effective support, while its ambient transport is the proper embedding
$\operatorname{Can}(m[S_\ell])\hookrightarrow S_a$.

\subsection{The Graph-Relative Canonicalizer}
\label{app:canon-g}

\leanref{TSG-APP-002}{TypedSlottedEGraphsPaper.ProfileRules.tsgApp002_completeCanonicalRecords}
Figure~\ref{alg:canon-g} gives the proposed graph-level pseudocode.  The exact
formal claim \textsc{app-002} concerns an abstract structural record extending
the extractor result with a shape and an effective-size renaming.  Its ambient
transport is definitionally the inclusion composed with that renaming, and its
collision key is definitionally the shape.  For any trace interpreter, replay
adds only equality of the two endpoint denotations.  No completeness of a
local orbit, graph insertion behavior, or $\operatorname{EqCert}$ construction
is part of \textsc{app-002}.  The paper tuple
\[
  (K_{\mathcal G}(n),\operatorname{Shape}_{\mathcal G}(n),\sigma_n,
   \iota_n,\omega_n,\xi_n)
\]
is the intended concrete realization of that record; establishing that the
Java record realizes it remains an artifact-refinement task.

\leanref{TSG-APP-003}{TypedSlottedEGraphsPaper.ProfileRules.tsgApp003_canonicalBlockOccurrenceContract}
The exact formal block-occurrence claim \textsc{app-003} uses de Bruijn
indices of one fixed finite size.  The canonical occurrence is the identity,
the alignment is the inverse of the source occurrence, and the returned list
is the descriptor's supplied automorphism list mapped by conjugation with that
identity.  Consequently source occurrence followed by alignment is the
identity occurrence.  Fresh-name choice, disjointness from a caller, leastness,
and completeness of any richer paper descriptor are not conclusions of this
claim; the corresponding steps in Figure~\ref{alg:canon-g} are refinement
requirements.

% Requires \usepackage{longtable,array,capt-of,needspace} in the document preamble.
\providecommand{\algline}[2][0pt]{%
  \hangindent=#1\hangafter=1\hspace*{#1}#2\strut\\%
}

\begingroup
\small
\setlength{\LTleft}{0pt}
\setlength{\LTright}{0pt}
\setlength{\LTpre}{0pt}
\setlength{\LTpost}{0pt}

% Reserve room for the caption, status strip, and opening procedure.
\Needspace{18\baselineskip}
\noindent\captionof{figure}{Proposed graph-relative canonicalization refinement
target (not verified as a producer by Lean).  The design minimizes invocation
and binder-block orbits before bag aggregation and set deduplication.  It would
retain $K_{\mathcal G}(n)$, $\iota_n$, $\omega_n$, and $\xi_n$ in the structural
record, obtain a dependent certificate $d_n^{\mathbf w}$ only by separate
replay, and use only $p_{\min}$ as a hash key.  The exact checked contracts are
stated after the figure.}
\label{alg:canon-g}

\begin{longtable}{@{}>{\raggedright\arraybackslash}p{0.98\linewidth}@{}}
\hline
\endfirsthead

\multicolumn{1}{@{}l@{}}{\small\itshape Figure~\thefigure\ continued}\\
\hline
\endhead

\hline
\endfoot

\endlastfoot

\algline{\textbf{status:} specification-only pseudocode; every producer and
  Java correspondence below is a separate refinement obligation}

\noalign{\vskip 0.35em}

\algline{\textbf{procedure} $\operatorname{leaderKernelTrace}_{\mathcal G}(n)$}
\algline[2.4em]{\textbf{require} $n$ is well typed and
  $\Gamma_0=\operatorname{slots}(n)$ is its exact context}
\algline[2.4em]{$(\widehat n_{\mathcal G},\xi)\leftarrow
  \operatorname{findAllWithTrace}_{\mathcal G}(n)$
  \hfill\textit{// find every invocation; do not unfold leaders}}
\algline[2.4em]{$(\widehat n_{\mathcal G},\xi)\leftarrow
  \operatorname{normalizeSignatureAdmittedPortsWithTrace}
  (\widehat n_{\mathcal G},\xi)$}
\algline[2.4em]{$\Delta\leftarrow\operatorname{slots}(\widehat n_{\mathcal G});
  \iota\leftarrow\operatorname{inclusion}(\Delta,\Gamma_0)$}
\algline[2.4em]{$n^\downarrow\leftarrow
  \operatorname{restrictExactContext}(\widehat n_{\mathcal G},\Delta)$}
\algline[2.4em]{\textbf{return} $(n^\downarrow,\iota,\xi)$}

\noalign{\vskip 0.35em}

\algline{\textbf{procedure}
  $\operatorname{canon}_{\mathcal G}(n=f_\theta(q_1,\ldots,q_r))$}
\algline[2.4em]{\textbf{require} $n$ is well typed and already uses its
  declared sibling quotients and explicit flat licenses}
\algline[2.4em]{$(K,\iota_n,\xi_n)\leftarrow
  \operatorname{leaderKernelTrace}_{\mathcal G}(n)$}
\algline[2.4em]{write $K=f_\theta(\bar q_1,\ldots,\bar q_r)$}
\algline[2.4em]{$\Delta\leftarrow\operatorname{slots}(K)$;
  $C\leftarrow\operatorname{Can}(\Delta)$}
\algline[2.4em]{$b\leftarrow\mathsf{None}$
  \hfill\textit{// option of a shape--renaming pair}}
\algline[2.4em]{\textbf{for each} $\sigma\in\operatorname{TRen}(C,\Delta)$ \textbf{do}}
\algline[4.8em]{$\rho\leftarrow\sigma^{-1}$
  \hfill\textit{// effective source slots to canonical slots}}
\algline[4.8em]{\textbf{for} $i\leftarrow1$ \textbf{to} $r$ \textbf{do}}
\algline[7.2em]{$r_i\leftarrow\operatorname{canonLeaderPort}_{\mathcal G}
  (\theta(\kappa_i),\bar q_i,\rho)$}
\algline[4.8em]{$p\leftarrow f_\theta(r_1,\ldots,r_r)$}
\algline[4.8em]{\textbf{if} $b=\mathsf{None}$ or
  $(p,\sigma)\mathrel{\triangleleft_{\mathrm{lex}}}\operatorname{get}(b)$
  \textbf{then} $b\leftarrow\mathsf{Some}(p,\sigma)$}
\algline[2.4em]{\textbf{require} $b=\mathsf{Some}(p_{\min},\sigma_{\min})$
  \hfill\textit{// $\operatorname{TRen}(C,\Delta)$ contains a bijection}}
\algline[2.4em]{$\omega_{\min}\leftarrow\iota_n\circ\sigma_{\min}$}
\algline[2.4em]{\textbf{return}
  $(K,p_{\min},\sigma_{\min},\iota_n,\omega_{\min},\xi_n)$}

\noalign{\vskip 0.35em}

\algline{\textbf{procedure} $\operatorname{replayKernelCertificate}_{\mathcal G,\mathbf w}
  (n,K,\iota,\xi)$}
\algline[2.4em]{\textbf{require} $\mathbf w$ is a coherent witness family for $\mathcal G$}
\algline[2.4em]{\textbf{require} $(K,\iota,\xi)=\operatorname{leaderKernelTrace}_{\mathcal G}(n)$
  and every certificate named by $\xi$ verifies}
\algline[2.4em]{$\Gamma_0\leftarrow\operatorname{slots}(n)$}
\algline[2.4em]{$d^{\mathbf w}\leftarrow
  \operatorname{replayTrace}_{\mathcal G,\mathbf w}(n,K,\iota,\xi)$}
\algline[2.4em]{\textbf{require}
  $d^{\mathbf w}:\operatorname{EqCert}_{\mathcal T_{\Sigma,\mathcal E}}
  (\Gamma_0;\lfloor n\rfloor_{\mathbf w},
  \iota\mathbin\cdot\lfloor K\rfloor_{\mathbf w})$}
\algline[2.4em]{\textbf{return} $d^{\mathbf w}$}

\noalign{\vskip 0.35em}

\algline{\textbf{procedure} $\operatorname{canon}_{\mathcal G}^{\mathbf w}(n)$}
\algline[2.4em]{$(K,p,\sigma,\iota,\omega,\xi)\leftarrow\operatorname{canon}_{\mathcal G}(n)$}
\algline[2.4em]{$d^{\mathbf w}\leftarrow
  \operatorname{replayKernelCertificate}_{\mathcal G,\mathbf w}
  (n,K,\iota,\xi)$}
\algline[2.4em]{\textbf{return}
  $(K,p,\sigma,\iota,\omega,\xi,d^{\mathbf w})$}

\noalign{\vskip 0.35em}

\algline{\textbf{procedure}
  $\operatorname{canonLeaderPort}_{\mathcal G}(\kappa,q,\rho)$}
\algline[2.4em]{\textbf{require}
  $q\in\operatorname{LPort}_{\mathcal G,\Lambda}(\kappa)$ and
  $\rho\in\operatorname{TRen}(\Lambda,\Lambda')$}
\algline[2.4em]{\textbf{target postcondition} result is
  $Q_{\mathcal G}^{\Lambda'}(\rho\mathbin\cdot q)$}
\algline[2.4em]{\textbf{match} $(\kappa,q)$ \textbf{with}}
\algline[4.8em]{$(\texttt{One}(\tau),\mathord{\$}x)$:\quad
  \textbf{return} $\rho(\mathord{\$}x)$}
\algline[4.8em]{$(\texttt{One}(\tau),m*\ell)$:}
\algline[7.2em]{\textbf{require} $\ell$ is a leader;
  $b\leftarrow\mathsf{None}$}
\algline[7.2em]{\textbf{for each} $\pi\in G_\ell$ \textbf{do}}
\algline[9.6em]{$u\leftarrow((\rho\circ m)\circ\pi)*\ell$}
\algline[9.6em]{$b\leftarrow\operatorname{keepLeast}_{\triangleleft}(b,u)$}
\algline[7.2em]{\textbf{return} $\operatorname{get}(b)$
  \hfill\textit{// $\operatorname{id}\in G_\ell$}}
\algline[4.8em]{$(\texttt{Seq}^{J}(\kappa'),[q_1,\ldots,q_k])$:}
\algline[7.2em]{\textbf{return}
  $[\operatorname{canonLeaderPort}_{\mathcal G}
    (\kappa',q_i,\rho)]_{i=1}^{k}$}
\algline[4.8em]{$(\texttt{Bag}^{J}(\kappa'),b)$:}
\algline[7.2em]{$b'\leftarrow$ the empty finite multiplicity map}
\algline[7.2em]{\textbf{for each} $q'\in\operatorname{supp}(b)$
  \textbf{do}}
\algline[9.6em]{$u\leftarrow\operatorname{canonLeaderPort}_{\mathcal G}
  (\kappa',q',\rho)$;
  $b'(u)\leftarrow b'(u)+b(q')$}
\algline[7.2em]{\textbf{return} $b'$
  \hfill\textit{// serialization order is not part of the bag}}
\algline[4.8em]{$(\texttt{Set}^{J}(\kappa'),Y)$:}
\algline[7.2em]{\textbf{return}
  $\{\operatorname{canonLeaderPort}_{\mathcal G}
    (\kappa',q',\rho)\mid q'\in Y\}$}
\algline[4.8em]{$(\texttt{Bind}(\tau,\kappa'),
  \operatorname{bind}\mathord{\$}x{:}\tau.q)$:}
\algline[7.2em]{$\mathord{\$}y\leftarrow
  \operatorname{freshCanonicalSlot}(\tau,\operatorname{cod}(\rho))$}
\algline[7.2em]{$\rho'\leftarrow
  \rho\mathbin{\uplus}[\mathord{\$}x\mapsto\mathord{\$}y]$}
\algline[7.2em]{\textbf{return}
  $\operatorname{bind}\mathord{\$}y{:}\tau.
  \operatorname{canonLeaderPort}_{\mathcal G}(\kappa',q,\rho')$}
\algline[4.8em]{$(\texttt{BindBlock}(\beta,\kappa'),
  \operatorname{bind}_{\beta}^{o}q)$:}
\algline[7.2em]{$(o_B,\alpha_o,A_B)\leftarrow
  \operatorname{canonicalBlockOccurrence}(\beta,o,
  \operatorname{cod}(\rho))$}
\algline[7.2em]{\hspace{1em}\textit{// $o_B=(B,a_B)$ is fresh;
  $\alpha_o=a_B\circ a_o^{-1}$; $A_B$ is the transported group}}
\algline[7.2em]{$b\leftarrow\mathsf{None}$}
\algline[7.2em]{\textbf{for each} $\pi_B\in A_B$ \textbf{do}}
\algline[9.6em]{$\rho_{\pi}\leftarrow
  \rho\mathbin{\uplus}(\pi_B\circ\alpha_o)$}
\algline[9.6em]{$u\leftarrow\operatorname{bind}_{\beta}^{o_B}
  \operatorname{canonLeaderPort}_{\mathcal G}(\kappa',q,\rho_{\pi})$}
\algline[9.6em]{$b\leftarrow\operatorname{keepLeast}_{\triangleleft}(b,u)$}
\algline[7.2em]{\textbf{return} $\operatorname{get}(b)$
  \hfill\textit{// $\operatorname{id}\in\operatorname{Aut}(\beta)$}}
\hline
\end{longtable}
\endgroup


\paragraph{Exact formal contracts for the figure.}
The algorithm rows are read through the following abstract claims, not as a
verified execution of the displayed graph program.
\leanref{TSG-ALG-005}{TypedSlottedEGraphsPaper.ProfileRules.tsgAlg005_leaderKernelTracePseudocodeContract}
\textsc{alg-005} says that an abstract kernel extractor's returned finite
inclusion is injective and that its trace replays to endpoint-denotation
equality under a supplied interpreter.
\leanref{TSG-ALG-006}{TypedSlottedEGraphsPaper.ClaimLedger.TSGALG006CompleteCanonPseudocodeContract}
\textsc{alg-006} uses an extensionally complete nonempty finite candidate list:
every enumerated candidate is admissible, every admissible candidate is
enumerated, the selected candidate is present and least, the returned record
contains exactly it, and the collision key projects only the returned shape.
These are finite-presentation contracts, not a proof that the concrete
$Q_{\mathcal G}$ enumerator has them.
\leanref{TSG-ALG-007}{TypedSlottedEGraphsPaper.ProfileRules.tsgAlg007_replayKernelCertificatePseudocodeContract}
\textsc{alg-007} says that a structural trace between two dependent sigma
endpoints replays to equality under a supplied interpreter.
\leanref{TSG-ALG-008}{TypedSlottedEGraphsPaper.ProfileRules.tsgAlg008_certifiedCanonWrapperPseudocodeContract}
\textsc{alg-008} says that the certified wrapper preserves its input
structural record and adds that replayed equality.
\leanref{TSG-ALG-009}{TypedSlottedEGraphsPaper.ProfileRules.tsgAlg009_canonLeaderPortPseudocodeContract}
Finally, \textsc{alg-009} proves constructor-by-constructor that the abstract
recursive function realizes its independently defined \emph{untyped}
$\operatorname{PortQuotient}$ relation: slots are mapped, invocation candidates
are passed to the defined $\operatorname{leastByKey}$ function, sequence and bag lists are recursively
mapped, set lists are deduplicated by key, unary binders recurse, and a block
applies the same function to its supplied candidates.  It does not prove typed-schema
admissibility, graph lookup, or completeness of the supplied invocation or
block candidates.

\leanref{TSG-APP-004}{TypedSlottedEGraphsPaper.ProfileRules.tsgApp004_canonLeaderPortContract}
The exact finite-orbit claim \textsc{app-004} is correspondingly generic.  For
any finite normalization presentation $P$, $P.\mathit{normal}(x)$ belongs to
the supplied orbit of $x$ and is $P$-related to $x$, and
\[
  P.\mathit{normal}(x)=P.\mathit{normal}(y)
  \quad\Longleftrightarrow\quad P.\mathit{relation}(x,y).
\]
This is the finite-presentation normal-form theorem used by the paper; identifying $P$ with the
concrete graph-relative relation is an additional refinement obligation.

\leanref{TSG-APP-005}{TypedSlottedEGraphsPaper.ClaimLedger.TSGAPP005CertifiedAbstractRebuildContract}
The exact rebuild claim \textsc{app-005} is an abstract three-part contract.
For a supplied finite-decreasing rebuild system, every initial state reaches
some quiescent state.  A supplied interface-restriction certificate decomposes
the old natural-number size as the retained size plus further removals plus
one, so committing it strictly decreases the size.  Finally, pointwise port
evidence at one shared head and type instantiation constructs the corresponding
node-congruence witness.  Concrete collision buckets, transport of graph
witnesses, pending-certificate behavior, and Java fixed-point scheduling are
not conclusions of this theorem; they remain implementation obligations.

\subsection{Ordered Application of Rules}
\label{app:rule-order}

\leanref{TSG-RULE-001}{TypedSlottedEGraphsPaper.ProfileRules.tsgRule001_orderedRulePipeline}
The formal object for \textsc{rule-001} is the displayed list of nine stage
names: (1) parser cleanup; (2) alpha normalization; (3) beta reduction;
(4) branch elimination; (5) first NNF; (6) phase-local prenexing; (7) guard
elimination and NNF; (8) sibling and flat normalization; and (9) local
saturation and rebuild.  Its only phase theorem is that the formal
$\operatorname{normalizeWithinPhase}$ operation, presently the identity,
preserves the stored phase index.  It does not establish that the Java stages
execute in this order, that temporal children are normalized recursively, or
that arbitrary quantifier motion cannot cross a phase.  Those are artifact
conformance requirements.

\subsection{Parser and Structural Normalization}

\leanref{TSG-RULE-002}{TypedSlottedEGraphsPaper.ProfileRules.tsgRule002_parserHeadNormalization}
For \textsc{rule-002}, each finite parser-head constructor in
Table~\ref{tab:opcode-aliases} maps to the shown canonical-head constructor;
the enumerated operation-meaning label of the result equals the enumerated
label of the source.  This is a datatype equation, not a denotational theorem
about the Java parser or Alloy expressions.

\begin{table}[t]
\centering
\small
\caption{Parser heads mapped to canonical internal heads.}
\label{tab:opcode-aliases}
\setlength{\tabcolsep}{4pt}
\begin{tabular}{ll@{\qquad}ll}
\hline
Source heads & Canonical head & Source heads & Canonical head\\
\hline
\texttt{BF/AND}, \texttt{LF/AND} & \texttt{BOOL/AND} &
\texttt{BF/OR}, \texttt{LF/OR} & \texttt{BOOL/OR}\\
\texttt{BF/IMPLIES} & \texttt{BOOL/IMPLIES} &
\texttt{BF/IFF} & \texttt{BOOL/IFF}\\
\texttt{UF/NOT} & \texttt{BOOL/NOT} &
\texttt{BE/PLUS} & \texttt{REL/PLUS}\\
\texttt{BE/INTERSECT} & \texttt{REL/INTERSECT} &
\texttt{BE/JOIN} & \texttt{REL/JOIN}\\
\texttt{BE/ARROW} & \texttt{REL/ARROW} &
\texttt{BE/MUL} & \texttt{ARITH/MUL}\\
\texttt{BE/IPLUS} & \texttt{ARITH/PLUS} &
\texttt{LE/DISJOINT} & \texttt{LIST/DISJOINT}\\
\hline
\end{tabular}
\end{table}

\leanref{TSG-RULE-003}{TypedSlottedEGraphsPaper.ProfileRules.tsgRule003_parserStructuralCleanup}
Internal no-op and end markers are removed:
\[
  \operatorname{NOOP}(E)\longrightarrow E,
  \qquad
  \langle\mathrm{END}\rangle\longrightarrow\epsilon.
\]
The exact \textsc{rule-003} equations say that a singleton \textsc{noop} list
unwraps, a singleton \textsc{end} list deletes, and a singleton abstract
\textsc{incomplete} list deletes.  They contain no arity validator and no claim
about Java parser rejection.  Lexical alpha-renaming and diagnostic-name
retention are likewise implementation requirements, not conclusions of this
rule.

\leanref{TSG-RULE-004}{TypedSlottedEGraphsPaper.ProfileRules.tsgRule004_captureAvoidingLetBeta}
For intrinsically scoped de Bruijn terms, \textsc{rule-004} defines
\[
  \operatorname{let}\ x=E\mid P\longrightarrow P[E/x].
\]
Precisely, reducing the formal $\mathsf{letE}(E,P)$ returns
$\mathsf{some}(\operatorname{substituteTop}(E,P))$, and lifted substitution
fixes the newly bound coordinate $0$.  The index types enforce scoping; the
theorem does not assert semantic preservation for a concrete Alloy AST or its
Java substitution routine.

\subsection{Branch Connectives and Boolean NNF}

\leanref{TSG-RULE-005}{TypedSlottedEGraphsPaper.ProfileRules.tsgRule005_branchConnectiveElimination}
The third rule below applies to formula-valued conditionals:
\begin{align*}
 A\Rightarrow B
   &\longrightarrow \neg A\lor B,\\
 A\Leftrightarrow B
   &\longrightarrow (\neg A\lor B)\land(\neg B\lor A),\\
 \operatorname{ite}(C,A,B)
   &\longrightarrow (C\land A)\lor(\neg C\land B).
\end{align*}
The same Boolean equations imply
\begin{align*}
 \neg(A\Rightarrow B)&\longrightarrow A\land\neg B,\\
 \neg(A\Leftrightarrow B)&\longrightarrow
   (A\land\neg B)\lor(B\land\neg A).
\end{align*}
These five \textsc{rule-005} equalities are proved only for the formal Boolean
truth functions.  Classifying and rewriting formula-valued versus
expression-valued Alloy conditionals is an artifact-refinement obligation.

\leanref{TSG-RULE-006}{TypedSlottedEGraphsPaper.ProfileRules.tsgRule006_booleanNegationNormalForm}
The Boolean NNF rules are
\begin{align*}
 \neg\top&\longrightarrow\bot,&
 \neg\bot&\longrightarrow\top,&
 \neg\neg A&\longrightarrow A,\\
 \neg\bigwedge_i A_i&\longrightarrow\bigvee_i\neg A_i,&
 \neg\bigvee_i A_i&\longrightarrow\bigwedge_i\neg A_i.
\end{align*}
Here \textsc{rule-006} is exactly the two constant equations, double negation,
and the two De Morgan equations for finite lists of Booleans.

\leanref{TSG-RULE-007}{TypedSlottedEGraphsPaper.ProfileRules.tsgRule007_involutiveAtomicNegationDuals}
Atomic negation uses the following involutive duals:
\[
\begin{array}{c@{\quad\leftrightarrow\quad}c@{\qquad}c@{\quad\leftrightarrow\quad}c}
 = & \ne & > & \le \\
 \ge & < & \in & \notin \\
 \multicolumn{4}{c}{\operatorname{some}R\quad\leftrightarrow\quad\operatorname{no}R}
\end{array}
\]
The exact \textsc{rule-007} theorem is conditional: if the finite
$\operatorname{atomicDual}$ function returns $o^\bot$ for $o$, it returns $o$
for $o^\bot$, and normalizing the negation of $o$ returns the atom
$o^\bot$.  Behavior of parser-specific negative comparison heads and
unsupported operators is not part of this theorem; Java conformance to the
table is checked separately.

\subsection{Temporal Negation}

\leanref{TSG-RULE-008}{TypedSlottedEGraphsPaper.ProfileRules.tsgRule008_temporalNegationDuals}
The supported temporal duals are
\begin{align*}
 \neg\operatorname{always}A
   &\longrightarrow\operatorname{eventually}\neg A,\\
 \neg\operatorname{eventually}A
   &\longrightarrow\operatorname{always}\neg A,\\
 \neg\operatorname{historically}A
   &\longrightarrow\operatorname{once}\neg A,\\
 \neg\operatorname{once}A
   &\longrightarrow\operatorname{historically}\neg A,\\
 \neg(A\ \mathsf{until}\ B)
   &\longrightarrow(\neg A)\ \mathsf{releases}\ (\neg B),\\
 \neg(A\ \mathsf{releases}\ B)
   &\longrightarrow(\neg A)\ \mathsf{until}\ (\neg B),\\
 \neg(A\ \mathsf{since}\ B)
   &\longrightarrow(\neg A)\ \mathsf{triggered}\ (\neg B),\\
 \neg(A\ \mathsf{triggered}\ B)
   &\longrightarrow(\neg A)\ \mathsf{since}\ (\neg B).
\end{align*}
The exact \textsc{rule-008} theorem has the same conditional form: whenever
$\operatorname{temporalDual}(o)=o^\bot$, the reverse lookup returns $o$ and
normalization returns the dualized constructor $o^\bot$.  It does not prove a
temporal semantics, phase separation, or Java handling of \textsf{before} and
\textsf{after}.

\subsection{Quantifier Rewrites and Prenexing}

\leanref{TSG-RULE-009}{TypedSlottedEGraphsPaper.ProfileRules.tsgRule009_quantifierNegation}
Write $\mathcal Qx{:}D.P$ suggestively, but \textsc{rule-009} is formalized as
evaluation of a quantifier over a finite list of Boolean matrix values.
Negation follows exactly
\begin{align*}
 \neg(\forall x{:}D.P)&\longrightarrow\exists x{:}D.\neg P,\\
 \neg(\exists x{:}D.P)&\longrightarrow\forall x{:}D.\neg P,\\
 \neg(\mathsf{no}\ x{:}D.P)&\longrightarrow\exists x{:}D.P,\\
 \neg(\mathsf{one}\ x{:}D.P)&\longrightarrow\mathsf{notone}\ x{:}D.P,\\
 \neg(\mathsf{notone}\ x{:}D.P)&\longrightarrow\mathsf{one}\ x{:}D.P,\\
 \neg(\mathsf{lone}\ x{:}D.P)&\longrightarrow\mathsf{notlone}\ x{:}D.P,\\
 \neg(\mathsf{notlone}\ x{:}D.P)&\longrightarrow\mathsf{lone}\ x{:}D.P.
\end{align*}
The theorem is the corresponding Boolean-list equality for all seven formal
quantifier constructors; it is not an Alloy evaluator correspondence theorem.

\leanref{TSG-RULE-010}{TypedSlottedEGraphsPaper.ProfileRules.tsgRule010_emptyAdmissibleBindings}
For an empty domain, independently of $P$,
\[
\begin{array}{rcl@{\qquad}rcl}
 \forall x{:}\varnothing.P&\longrightarrow&\top &
 \exists x{:}\varnothing.P&\longrightarrow&\bot\\
 \mathsf{no}\ x{:}\varnothing.P&\longrightarrow&\top &
 \mathsf{one}\ x{:}\varnothing.P&\longrightarrow&\bot\\
 \mathsf{lone}\ x{:}\varnothing.P&\longrightarrow&\top &
 \mathsf{notone}\ x{:}\varnothing.P&\longrightarrow&\top\\
 \mathsf{notlone}\ x{:}\varnothing.P&\longrightarrow&\bot.&&
\end{array}
\]
For \textsc{rule-010}, ``empty'' means exactly that the formal admissible-value
list is $[]$; the seven displayed values are then proved by evaluation.
Determining emptiness from Alloy declaration syntax and multiplicity is outside
this theorem and is a required frontend/Java check.

\leanref{TSG-RULE-011}{TypedSlottedEGraphsPaper.ProfileRules.tsgRule011_constantBodyQuantifiers}
The carrier-independent constant-body rules are
\[
\begin{array}{rcl@{\qquad}rcl}
 \forall x{:}D.\top&\longrightarrow&\top &
 \exists x{:}D.\bot&\longrightarrow&\bot\\
 \mathsf{no}\ x{:}D.\bot&\longrightarrow&\top &
 \mathsf{one}\ x{:}D.\bot&\longrightarrow&\bot\\
 \mathsf{lone}\ x{:}D.\bot&\longrightarrow&\top &
 \mathsf{notone}\ x{:}D.\bot&\longrightarrow&\top\\
 \mathsf{notlone}\ x{:}D.\bot&\longrightarrow&\bot.&&
\end{array}
\]
These are exactly the seven \textsc{rule-011} equalities for a list of any
length.  No value is asserted for existential true or universal false, whose
Boolean-list values depend on whether that list is empty.

\leanref{TSG-RULE-012}{TypedSlottedEGraphsPaper.ProfileRules.tsgRule012_phaseLocalPrenexAndDomainGuarding}
The exact \textsc{rule-012} statement is four finite Boolean-list equations.
For a carrier list $C$, predicates $p,d:C\to\mathbb B$, and Boolean $r$,
\begin{align*}
 (\exists_C p)\land r&=\exists_C(\lambda x.\,p(x)\land r),&
 (\forall_C p)\lor r&=\forall_C(\lambda x.\,p(x)\lor r),\\
 \exists_{\operatorname{filter}(d,C)}p
   &=\exists_C(\lambda x.\,d(x)\land p(x)),&
 \forall_{\operatorname{filter}(d,C)}p
   &=\forall_C(\lambda x.\,\neg d(x)\lor p(x)).
\end{align*}
Its binding-tuple operation changes only the primitive carrier and preserves
the quantifier, lower and upper multiplicities, disjointness class, and phase
fields.  Fresh-variable conditions, general AST free-variable preservation,
cardinality lowering, sums, and comprehensions are not conclusions of this
claim.

\leanref{TSG-RULE-013}{TypedSlottedEGraphsPaper.ProfileRules.tsgRule013_descriptorCertifiedBlockPermutation}
For \textsc{rule-013}, an accepted finite permutation is packaged with a proof
that the supplied binder descriptor certifies it.  The theorem proves that it
occurs in the descriptor's supplied automorphism list and that, for every
environment coordinate $i$,
\[
  \mathit{env}\bigl(\pi^{-1}(\pi(i))\bigr)=\mathit{env}(i).
\]
It does not establish semantic preservation of a quantified Alloy body or
preservation of additional descriptor metadata; those are separate refinement
obligations.

\subsection{Flexible-Arity Structural Laws}

\leanref{TSG-RULE-014}{TypedSlottedEGraphsPaper.ProfileRules.tsgRule014_completeOperatorLawBoundary}
For \textsc{rule-014}, Table~\ref{tab:operator-laws} is exactly the finite
$\operatorname{operatorLaw}$ lookup table.  Its three Boolean columns are
policy flags, not semantic certificates.  In particular, the theorem proves
the noninteger lookup equations; the integer row belongs to
\textsc{rule-015}.  Neither theorem derives a source-level law merely from a
flag.

\begin{table}[t]
\centering
\footnotesize
\caption{Formal carrier assignments and policy flags.}
\label{tab:operator-laws}
\setlength{\tabcolsep}{4pt}
\begin{tabular}{p{0.25\linewidth}p{0.20\linewidth}p{0.22\linewidth}p{0.17\linewidth}}
\hline
Formal operator & Carrier & C/I flags & Flat flag\\
\hline
$\mathsf{boolAnd},\mathsf{boolOr}$ & $\mathsf{setPlus}$ & true/true & true\\
$\mathsf{relationalUnion}$,\newline $\mathsf{relationalIntersection}$ & $\mathsf{setPlus}$ & true/true & true\\
$\mathsf{integerAddition}$,\newline $\mathsf{integerMultiplication}$ & $\mathsf{bagPlus}$ & true/false & profile-dependent\\
$\mathsf{relationalJoin}$,\newline $\mathsf{relationalProduct}$ & $\mathsf{fixedBinary}\allowbreak\mathsf{Positional}$ & false/false & false\\
$\mathsf{equality},\mathsf{inequality}$,\newline $\mathsf{biconditional}$ & $\mathsf{fixedBinary}\allowbreak\mathsf{Commutative}$ & true/false & false\\
$\mathsf{disjoint}(k)$ & $\mathsf{nonflatBag}$ & true/false & false\\
$\mathsf{call}(g,k)$ & $\mathsf{positional}(k)$ & false/false & false\\
$\mathsf{totalOrder}(k)$ & $\mathsf{positional}(k)$ & false/false & false\\
\hline
\end{tabular}
\end{table}

The disjoint, call, and total-order equations are quantified over their natural
number arities (and, for calls, the callee identifier).  Whether the Java
frontend enforces source arities and whether a policy flag is backed by the
required endpoint certificates are separate artifact checks.  No additional
role-sensitive-list or dependent-chain row is covered by \textsc{rule-014}.

\leanref{TSG-RULE-015}{TypedSlottedEGraphsPaper.ProfileRules.tsgRule015_modularIntegerACProfile}
The exact \textsc{rule-015} model uses $\operatorname{Fin}(2^w)$.  Its addition
and multiplication are proved associative and commutative; the modular profile
has both flat flags set, and the overflow-forbidding profile has both cleared.
\leanref{TSG-CE-003}{TypedSlottedEGraphsPaper.ProfileRules.tsgCe003_overflowReassociationCounterexample}
Separately, the checked signed four-bit functions prove
$\operatorname{checkedAdd4}(7,1)=\mathsf{none}$,
$\operatorname{checkedAdd4}(1,-1)=\mathsf{some}(0)$, and hence the left
association fails while the right returns $7$.  Connecting either model to an
Alloy or Java integer mode is not part of these Lean results.

In the abstract \textsc{alg-009} procedure, sequence and bag constructors map
their child lists recursively, whereas a set constructor additionally applies
the supplied key-deduplication function.  This statement is about lists in the
formal quotient datatype; it does not prove multiset semantics, Java sorting,
or graph-invocation equality.

\subsection{Local Saturation Rules}

\leanref{TSG-RULE-016}{TypedSlottedEGraphsPaper.ProfileRules.tsgRule016_booleanLocalSaturation}
The Boolean identities, annihilators, and complements are
\begin{align*}
 A\land A&\longrightarrow A,& A\lor A&\longrightarrow A,\\
 A\land\top&\longrightarrow A,& A\lor\bot&\longrightarrow A,\\
 A\land\bot&\longrightarrow\bot,& A\lor\top&\longrightarrow\top,\\
 A\land\neg A&\longrightarrow\bot,&
 A\lor\neg A&\longrightarrow\top.
\end{align*}
These eight \textsc{rule-016} identities are Boolean equations.  The associated
formal smart constructor returns a constant on an empty input list, the child
on a singleton, and a nonempty-port record otherwise; its stored-port arity can
never be $0$.  No theorem here connects that datatype to a frontend
$\texttt{Set}^+$ node or to an emitted unit certificate.

\leanref{TSG-RULE-017}{TypedSlottedEGraphsPaper.ProfileRules.tsgRule017_atomicComplementSaturation}
For \textsc{rule-017}, an atomic literal consists only of an arbitrary atom and
a positive/negative polarity.  Complement flips the polarity while retaining
the same atom, and for any Boolean valuation
\[
  \ell\land\overline\ell=\bot,
  \qquad
  \ell\lor\overline\ell=\top.
\]
Recognizing concrete operator duals and equal canonical e-class invocations is
an implementation premise, not a conclusion of this theorem.

\leanref{TSG-RULE-018}{TypedSlottedEGraphsPaper.ProfileRules.tsgRule018_relationalLocalSaturation}
For \textsc{rule-018}, relations are abstract predicates
$\alpha\to\mathsf{Prop}$.  Given a witness that $x$ is nonempty, the theorem
proves pointwise
\begin{align*}
 x\in\mathsf{none}&\longrightarrow\bot,&
 x\in\mathsf{univ}&\longrightarrow\top,\\
 R+\mathsf{none}&\longrightarrow R,&
 R\mathbin{\&}\mathsf{none}&\longrightarrow\mathsf{none},\\
 R+R&\longrightarrow R,&
 R\mathbin{\&}R&\longrightarrow R.
\end{align*}
and also that the empty predicate is a subset of itself.  This theorem contains
no typed Alloy AST, multiplicity analysis, ACI-port inference, or Java rewrite
claim.

\leanref{TSG-RULE-019}{TypedSlottedEGraphsPaper.ProfileRules.tsgRule019_implicationLocalSaturation}
Finally, \textsc{rule-019} proves the Boolean identities
\begin{align*}
 \bot\Rightarrow A&\longrightarrow\top,&
 \top\Rightarrow A&\longrightarrow A,\\
 A\Rightarrow\top&\longrightarrow\top,&
 A\Rightarrow\bot&\longrightarrow\neg A,
\end{align*}
and the Boolean equality $A\Rightarrow B=\neg A\lor B$.  Stage placement and
Java saturation behavior are outside this equality theorem.

\subsection{Boundary and Implementation Correspondence}

\leanref{TSG-SCOPE-001}{TypedSlottedEGraphsPaper.ProfileRules.tsgScope001_zeroCanonicalDistanceBoundary}
The exact \textsc{scope-001} theorem is about an arbitrary complete finite
normalization presentation $P$.  Define its discrete canonical distance to be
$0$ when $P.\mathit{normal}(x)=P.\mathit{normal}(y)$ and $1$ otherwise.  Then
\[
  \operatorname{canonicalDistance}_P(x,y)=0
  \quad\Longleftrightarrow\quad P.\mathit{relation}(x,y).
\]
This follows from the presentation's normal-form exactness.  It neither proves
that the individual rules in this appendix generate one combined $P$, nor
identifies this $0/1$ function with the dataset's AST edit distance, nor states
a Java-refinement premise or conclusion.

The remaining production statements are artifact-only obligations.  In
particular, graph-relative typed-renaming enumeration, binder-group closure,
the behavior of \texttt{SlotPermutationGroup}, the 720-permutation guard, and
the 32-pass rewrite bound are not Lean conclusions.  The publication run gives
in-process checked observations under its recorded rule set; it does not prove
those correspondences and cannot be cited as realizing \textsc{scope-001} or
the abstract algorithm contracts above.
